\documentclass[a4paper,11pt]{article}
\usepackage[utf8]{inputenc}
\usepackage[T1]{fontenc}
\usepackage[french]{babel}
\usepackage{geometry}
\geometry{left=1.5cm,right=1.5cm,top=1.2cm,bottom=1.2cm}

\usepackage{lmodern}
\usepackage{xcolor}
\definecolor{titleblue}{RGB}{0,51,140}
\definecolor{TITLEBLUE}{RGB}{0,51,140}

\usepackage{hyperref}
\hypersetup{colorlinks=true, urlcolor=titleblue, linkcolor=titleblue}

\usepackage{titlesec}
\usepackage{enumitem}

\pagestyle{empty}

\titleformat{\section}{\color{titleblue}\Large\bfseries\uppercase}{}{0em}{}[\color{titleblue}\titlerule]
\titlespacing\section{0pt}{8pt}{4pt}
\setlist[itemize]{leftmargin=*,topsep=1pt,itemsep=1pt}

\begin{document}

\begin{center}
    {\Huge\bfseries\color{titleblue} AISSAOUI Ismail}\\[4pt]
    {\Large Ingénieur d'État en Intelligence Artificielle \& Data Science}\\[6pt]
    ismail.aissaoui.pro@gmail.com $|$ +213 660 70 77 96 $|$ Chlef, Algérie \\
    \href{https://linkedin.com/in/aissaoui-ismail-6a77a92a7}{linkedin.com/in/aissaoui-ismail-6a77a92a7} $|$ \href{https://github.com/i-aissaoui}{github.com/i-aissaoui}
\end{center}

\section{RÉSUMÉ PROFESSIONNEL}
Ingénieur diplômé de l'ESI-SBA, spécialisé dans les \textbf{Jumeaux Numériques}, la visualisation 3D et l'IA temps réel. Expertise en optimisation de la latence "Edge-to-User" et architectures distribuées pour systèmes cyber-physiques (CPS). Candidat au doctorat focalisé sur la compression sémantique pour la XR industrielle.

\section{FORMATION}
\textbf{École Nationale Supérieure d’Informatique (ESI-SBA)} \hfill 2021 -- 2026 \\
\textbf{Diplôme d'Ingénieur d'État \& Master en IA}

\section{EXPÉRIENCE PROFESSIONNELLE}
\textbf{Stagiaire Ingénieur IA — Jumeau Numérique 3D \& Edge Computing} \hfill Jan 2026 -- Présent \\
\textbf{Sonatrach (Branche Transport par Canalisations)} \\
\begin{itemize}
    \item Conception d'une interface de monitoring 3D (\textbf{Unity/Three.js context}) synchronisée.
    \item Optimisation de la latence d'inférence à \textbf{0.04ms} via \textbf{C++ JIT xLSTM}.
    \item Développement de substituts génératifs (\textbf{Mamba-SSM}) pour simulateurs temps réel.
    \item Co-conception Matériel-Logiciel pour périphériques mobiles contraints.
\end{itemize}

\textbf{Stagiaire en ingénierie réseau} \hfill Sep 2024 \\
\textbf{Algérie Télécom RMC Center}

\section{PROJETS DE RECHERCHE ET INGÉNIERIE}

\textbf{Geo-RAG — Geospatial Knowledge \& Reasoning} \hfill 2025 \\
Système intelligent utilisant le RAG pour le traitement de données sémantiques multi-sources.
\textit{Compétence : Analyse sémantique et extraction de connaissances.}

\textbf{Recognizini — Reconnaissance faciale \& Vision Temps Réel} \hfill 2025 \\
Système d'apprentissage actif avec embeddings incrémentaux pour la vision par ordinateur.

\textbf{Career Connect — GNN recruitment Intelligence} \hfill 2024 -- 2025 \\
Modélisation de relations complexes via Graph Neural Networks.

\textbf{University Scheduling (NP-Complete Optimization)} \hfill 2025 \\
Optimisation de systèmes complexes via métaheuristiques (\textbf{PSO/ACO/GA}).

\textbf{Aura — Reinforcement Learning \& Distributed Systems} \hfill 2025 \\
Plateforme distribuée avec dashboard temps réel et agents autonomes.

\section{COMPÉTENCES TECHNIQUES}
\begin{itemize}
    \item \textbf{XR \& Graphics}: Unity 3D, Three.js, Electron, Visualisation de données 3D.
    \item \textbf{IA Temps Réel}: Mamba-SSM, xLSTM, Inférence Edge (C++), Optimisation de latence.
    \item \textbf{Data Science}: Apprentissage Sémantique, Compression, GNN, PINN, RL.
    \item \textbf{Stack}: Python (PyTorch), C++, Docker, Node.js, Git, Linux.
\end{itemize}

\section{LANGUES}
Arabe (Maternel), Anglais (Avancé), Français (Courant).

\end{document}
